\section{Exact parameters of concrete affine formulations}
\label{sec:concrete-barriers}

This section computes barrier parameters for three formulations of
balls, each on its own affine domain: norm trees, grouped slices, and
column packing.  The norm tree uses the fewest small cone factors, but
its restricted standard barrier often has a larger parameter than that
of the grouped formulation.  For the grouped and column-packed
slices we also prove lower bounds for every self-concordant barrier on
the same domain, including barriers with arbitrary mixed derivatives.
For an open convex domain $\Omega$ in its affine hull, the optimal, or
intrinsic, parameter $\vartheta_{\mathrm{opt}}(\Omega)$ is the infimum of
the parameters of all self-concordant barriers with positive definite
Hessian on that affine hull; the formulation stays fixed.

These standard-parameter calculations also appear in the unpublished companion manuscript
\emph{The Cost of Following the Central Path}~\citep{CentralPathCompanion}.
We reproduce them with proofs so that this paper is self-contained, and
use them to study arbitrary barriers on a fixed affine domain and to
compare resources across lifts.  We do not claim the shared
calculations and constructions as new.

\subsection{Norm trees: the children of the root determine the standard parameter}

A norm tree for $\B_2^N$ uses a rooted tree $T$ with $N\geq2$ leaves,
one per coordinate, and $L$ internal nodes, each with $c_v\geq2$
children and one Lorentz cone.  At an internal node $v$, let $t_v$ be its axial
variable and let $y_v$ collect the values of its children: a leaf
contributes a projected coordinate, and an internal child $w$ contributes
$t_w$.  The linked cone is defined by
$t_v>\norm{y_v}$ for every internal node $v$.
Its root slice, with $t_r=1$, is denoted by $\Omega_T$ and projects onto
$\intr\B_2^N$.  Put
\[
 q_v=t_v^2-\norm{y_v}^2,\qquad
 \Phi_T=-\sum_v\log q_v,\qquad F_T=\Phi_T|_{t_r=1}.
\]
The positive-sheet condition $t_v>0$ is part of the domain.

\begin{theorem}[Standard parameter of a norm tree]\label{thm:tree-parameter}
The least self-concordance parameter of the restricted standard barrier
$F_T$ is
\begin{equation}\label{eq:tree-parameter-concrete}
 \nu(F_T)=
 \begin{cases}
  2L-2,&\text{every root child is internal},\\
  2L-1,&\text{at least one root child is a leaf}.
 \end{cases}
\end{equation}
For a forest of independent root-fixed trees, these parameters add.
\end{theorem}
\begin{proof}
Affine restriction of the standard Lorentz barriers gives standard
self-concordance.  The Hessian is positive definite, since every
nonzero linked direction changes a cone coordinate.  Before fixing
$t=t_r$, $\Phi_T$ is logarithmically homogeneous of degree $2L$.  With
$H=D^2\Phi_T$ at the point $z$ and $e_t$ the root covector, homogeneity
gives
$Hz=-D\Phi_T$ and $\norm{D\Phi_T}_{H^{-1}}^2=2L$.
Orthogonal projection onto the slice tangent space yields
\begin{equation}\label{eq:tree-root-leverage}
 \norm{DF_T}_{D^2F_T,*}^2
 =2L-\frac{t^2}{e_t^TH^{-1}e_t},\qquad
 \nu(F_T)=2L-\frac1{\sup\beta_T},\quad
 \beta_T=\frac{e_t^TH^{-1}e_t}{t^2}.
\end{equation}
The notation on the left means the dual norm induced by the restricted
Hessian.  The unit Dikin ellipsoid is contained in the linked cone,
which lies in $t>0$.  Consequently $e_t^TH^{-1}e_t\leq t^2$:
otherwise a vector of local norm less than one could move $t$ to zero.

If a root child is a leaf, set $t=1$, that leaf to $\sqrt{1-\epsilon}$,
and the other root leaves to zero.  Scale every internal child subtree
from a fixed interior point by $\epsilon^2$.  Then
$q_r=\epsilon-O(\epsilon^4)>0$.  The direction with $h_t=1$,
$h_{\mathrm{leaf}}=1$, and all other components zero has
\[
 D^2\Phi_T[h,h]
 =\frac{4(1-\sqrt{1-\epsilon})^2}{q_r^2}\longrightarrow1.
\]
The descendant terms vanish in this direction.  Since
$(e_t^TH^{-1}e_t)^{-1}=\min_{h_t=1}h^THh$, the preceding limit
and the Dikin bound prove $\sup\beta_T=1$.

Suppose instead that every root child is internal, with axial values
$a_i>0$.  Eliminate its strict descendants from the Hessian.  The
effective axial curvature of child $i$ is at least $a_i^{-2}$ by
Dikin containment in that child's linked cone.  Increasing these
curvatures can only increase the final root Schur complement.
It therefore suffices to consider
\[
 -\log\bigl(t^2-\textstyle\sum_i a_i^2\bigr)-\sum_i\log a_i.
\]
Normalize $t=1$, and put $r=\sum_i a_i^2$, $q=1-r$, and
$Z=\sum_i a_i^4/(q+2a_i^2)$.  The child Hessian is the diagonal
matrix with entries $2/q+a_i^{-2}$ plus $4aa^T/q^2$.
Its cross row with the root is $-4a^T/q^2$.  The rank-one inverse
formula gives the root Schur complement
\begin{equation}\label{eq:tree-schur-concrete}
 S=\frac{2(1+r)-16Z/(q+4Z)}{q^2}.
\end{equation}
Since $x/(q+2x)$ increases for $x\geq0$,
$Z\leq r^2/(q+2r)=r^2/(1+r)$.
The inequality $S\geq2$ is equivalent to
$r(3-r)\geq4(2-r)Z$; substituting this upper bound leaves
$3r(1-r)^2/(1+r)\geq0$.  Hence $\beta_T\leq1/2$.
Conversely, scale all child subtrees by $\epsilon\downarrow0$.
The direction changing only the root has squared local norm
$2(1+r)/(1-r)^2\to2$.  The same variational formula proves
$\sup\beta_T=1/2$.  Equation~\eqref{eq:tree-root-leverage} now proves
both cases.  On a Cartesian product, squared dual gradient norms add;
their suprema can be approached independently, proving the forest claim.
\end{proof}

\begin{corollary}[The best capped norm tree]\label{cor:tree-cap-parameter}
Fix an integer Lorentz dimension cap $d\geq3$, and set
\[
 A=d-2,\qquad E=N-1,\qquad L_0=\left\lceil E/A\right\rceil,
 \qquad \sigma=L_0A-E.
\]
Define $\chi=1$ if $L_0\geq3$ and
$\sigma\geq\max\{0,d-L_0\}$, and define $\chi=0$ otherwise.
Among all norm trees with $c_v+1\leq d$,
\begin{equation}\label{eq:tree-cap-parameter}
 \min_T\nu(F_T)=2L_0-1-\chi.
\end{equation}
\end{corollary}
\begin{proof}
The tree identity $\sum_v(c_v-1)=E$, with $1\leq c_v-1\leq A$,
gives $L\geq L_0$.  Any list of $L_0$ increments in this interval
with sum $E$ is realized by a chain of internal nodes, as in
Corollary~\ref{cor:dimension-resources}.
For an all-internal root of arity $c$, one must additionally have
$2\leq c\leq\min(d-1,L_0-1)$.  The remaining $L_0-1$ increments
must sum to $E-(c-1)$.  Such increments exist exactly when
\[
 L_0\geq3,\qquad
 E\leq (L_0-1)A+\min(d-1,L_0-1)-1.
\]
The lower bound on their sum causes no further constraint, since
$E\geq L_0$ and the root increment may be chosen no larger than
$E-L_0+1$.  Partitioning the remaining internal nodes into $c$
nonempty chains realizes the required descendant forest.
The displayed inequality is precisely the condition defining $\chi$.
Trees with more than $L_0$ internal nodes cannot do better: even with
an all-internal root, $L_0+1$ internal nodes give
$2(L_0+1)-2>2L_0-1$.
\end{proof}

The next proposition bounds every self-concordant barrier on the
root-fixed domain, not only the standard one.

\begin{proposition}[Intrinsic norm-tree bounds]\label{prop:tree-intrinsic}
Every self-concordant barrier on $\Omega_T$ has parameter at least $L$.
Every logarithmically homogeneous self-concordant barrier on the
unfixed linked cone has degree at least $L+1$.  For a forest of
independent root-fixed trees, every self-concordant barrier has
parameter at least the total number of internal nodes.
\end{proposition}
\begin{proof}
Choose positive leaf values with squared sum one and set $t_v^*$
equal to the norm of the leaves below $v$.  Each
$z_v^*=(t_v^*,y_v^*)$ is a nonzero boundary vector.  Intersect the
linked affine slice with the product of two-planes
\[
 z_v=\alpha_v(t_v^*,y_v^*)+\beta_v(t_v^*,-y_v^*).
\]
Cone membership is exactly $\alpha_v,\beta_v\geq0$; its determinant
is $4(t_v^*)^2\alpha_v\beta_v$.  Root fixing and linking become
\begin{equation}\label{eq:tree-polysection}
 \alpha_r+\beta_r=1,\qquad
 \alpha_w+\beta_w=\alpha_v-\beta_v
 \quad(w\text{ an internal child of }v).
\end{equation}
The $L$ variables $\beta_v$ are affine coordinates: these equations
determine every $\alpha_v$ recursively.  Small positive values of
all $\beta_v$ give strictly positive $\alpha_v$, so the section has
dimension $L$ and meets the interior of the linked slice.
Nonnegativity and the equations bound all coordinates by one.
At $\beta=0$, all $\alpha_v=1$, and exactly the $L$ independent
inequalities $\beta_v\geq0$ are active.  Every relative boundary
point of this polytope lies on a cone boundary.  Restriction of any
barrier on $\Omega_T$ is therefore a barrier on this polytope.
The facet bound \cite[Proposition~2.3.6]{NN1994} at this vertex, where
$L$ linearly independent facet inequalities are active, proves the
claim.  Product sections
prove the forest extension.

The corresponding independent-facet principle for logarithmically
homogeneous barriers on polyhedral cones, with one additional parameter
unit at a nonzero boundary point, is due to
\citet[Theorem~6.1]{Hildebrand2013}. The tree section above gives the
following direct argument for this application.
For the conic claim let the logarithmic degree be $\theta$.
Lemma~\ref{lem:compact-base-parameter} gives a root-slice parameter
at most $\theta-1$.  The first claim gives $\theta-1\geq L$.
Logarithmic homogeneity is essential to this last argument.
\end{proof}

Thus $\vartheta_{\mathrm{opt}}(\Omega_T)$ lies
between $L$ and the value in \eqref{eq:tree-parameter-concrete}.  We do
not determine it for general trees; in particular, we do not claim that
the standard barrier is optimal.

A natural attempt to lower the standard parameter adds $\log t_w$ for
every non-root internal node $w$, removing one axial logarithm per link.
This already fails for the smallest such tree, a root with a leaf child
$x$ and an internal child $a$ whose children are leaves $y$ and $z$.
The function
\[
 \widetilde F(t,x,a,y,z)
 =-\log(t^2-x^2-a^2)-\log(a^2-y^2-z^2)+\log a
\]
has logarithmic degree three.  At the interior point
$(3,0,1,0,0)$, in direction $(0,0,1,1,0)$, its derivatives are
$D^2\widetilde F=53/16$ and $D^3\widetilde F=-441/32$.
The squared self-concordance inequality fails, since
$(441/32)^2-4(53/16)^3=45604/1024>0$.
Hence this modification does not give a self-concordant barrier.

\subsection{Grouped fixed-tail slices: coupled barriers cannot lower the parameter}

For each ball $a$, partition its coordinates into nonempty groups
$G\in\mathcal G_a$, and put $H=\sum_a|\mathcal G_a|$.
Consider
\begin{equation}\label{eq:grouped-domain-concrete}
 \Omega_{\mathcal G}
 =\left\{(u,w):u_{aG}>\norm{w_{aG}}^2,\quad
                   \sum_{G\in\mathcal G_a}u_{aG}=1\text{ for every }a\right\}.
\end{equation}
Its projection is the product of the open balls.  The Lorentz realization
uses $((1+u)/2,(1-u)/2,w)\in\intr Q_{|G|+2}$.
The real PSD realization uses
$\left(\begin{smallmatrix}u&w^T\\w&I\end{smallmatrix}\right)\succ0$.
More generally, fixing the tail $z=1$ in
Lemma~\ref{lem:perspective} gives the same domain in every simple
Jordan family, for group size at most $a(r-1)$.
In all these cases the restricted determinant is
$u-\norm{w}^2$.

\begin{theorem}[Intrinsic grouped parameter]\label{thm:grouped-intrinsic}
The domain in \eqref{eq:grouped-domain-concrete} satisfies
\begin{equation}\label{eq:grouped-intrinsic}
 \vartheta_{\mathrm{opt}}(\Omega_{\mathcal G})=H.
\end{equation}
The restricted standard barrier
$F_{\mathcal G}=-\sum_{a,G}\log(u_{aG}-\norm{w_{aG}}^2)$
attains this value.
\end{theorem}
\begin{proof}
For $f(u,w)=-\log q$, $q=u-\norm{w}^2$, consider an affine line
and write $A=q'/q$, $C=2\norm{w'}^2/q\geq0$.  Then
\[
 f''=A^2+C,\quad f'''=-2A^3-3AC,\quad
 4(A^2+C)^3-(2A^3+3AC)^2=3A^2C^2+4C^3\geq0.
\]
The Hessian is positive definite, and
$(D^2f)^{-1}Df=(-q,0)$ gives squared dual gradient norm one.
Thus each summand is a one-self-concordant barrier; addition and
affine restriction give the upper bound $H$.

For the lower bound choose positive $\sigma_{aG}$ summing to one
for each $a$, and a unit vector $e_{aG}$ in each nonempty group.
The affine section $u_{aG}=\sigma_{aG}$,
$w_{aG}=t_{aG}e_{aG}$ is exactly the open box
\[
 \prod_{a,G}(-\sqrt{\sigma_{aG}},\sqrt{\sigma_{aG}}).
\]
Its boundary lies on the boundary of $\Omega_{\mathcal G}$.
An arbitrary barrier, including one coupling all groups, restricts
to a barrier on this $H$-dimensional box.  The classical facet
bound \cite[Proposition~2.3.6]{NN1994} yields $\nu\geq H$.
\end{proof}

For comparison, the product of $k$ positive-dimensional balls itself
has $\vartheta_{\mathrm{opt}}=k$.  The upper bound follows by restricting
the one-parameter paraboloid barriers to $u=1$ and adding; the lower
bound follows from a $k$-dimensional box section.  Consequently the
optimal parameter of the grouped domain exceeds that of its projection
by $H-k$.
Under a Lorentz dimension cap $d$ with $3\leq d<N+1$, the grouped
formulation of a single ball in $\R^N$ uses $h=\lceil N/(d-2)\rceil$
groups.  In the notation of
Corollary~\ref{cor:tree-cap-parameter}, put
$\zeta=\mathbf1_{\{d-2\mid N-1\}}$, so $h=L_0+\zeta$.
The exact comparison is
\begin{equation}\label{eq:tree-grouped-gap}
 \min_T\nu(F_T)-\nu(F_{\mathcal G})
 =L_0-1-\chi-\zeta\geq0.
\end{equation}
Here $L_0\geq2$; when $L_0=2$, $\chi=0$, and when $L_0\geq3$
the last inequality is immediate.  The grouped lift uses at most one
more factor than the smallest capped tree, while its parameter can be
asymptotically half the tree's.  If $d\geq N+1$, the direct Lorentz
slice already has parameter one.

The same box section also gives the lower bound $H$ for
$u_{aG}>\norm{w_{aG}}_p^p$, $1<p<\infty$, with the same allocation
equations.  The matching upper bound just proved uses the quadratic
case $p=2$; it is not asserted for other exponents.

\subsection{Bounded-fiber projection and one-block column packing}

A lower bound on the barrier parameter of a projected body implies one
for the lift only under hypotheses on the projection; counting lift
variables is not enough.  The next lemma, a form of the classical exact
partial-minimization theorem \cite[Theorem~5.2.1]{Chares2009}, gives
such hypotheses.

\begin{lemma}[Parameter monotonicity under bounded-fiber projection]
\label{lem:bounded-fiber-barrier}
Let $Q\subset E\times V$ be a closed convex set with nonempty interior,
and suppose every nonempty fiber $Q_x=\{y:(x,y)\in Q\}$ is bounded.
Suppose $C=\pi_E Q$ contains no line.  If $F$ is a self-concordant
barrier on $\intr Q$ with positive definite Hessian and parameter $\nu$,
then
\[
 \phi(x)=\min\{F(x,y):(x,y)\in\intr Q\}
\]
is a self-concordant barrier with positive definite Hessian and
parameter at most $\nu$ on $\intr C$.  The same statement holds in
affine hulls after eliminating affine equations.  In particular,
$\vartheta_{\mathrm{opt}}(\intr Q)
\geq\vartheta_{\mathrm{opt}}(\intr C)$.
\end{lemma}
\begin{proof}
First, inverse images in $Q$ of bounded subsets of $E$ are bounded.
Otherwise choose $(x_j,y_j)\in Q$ with $x_j$ bounded and
$\norm{y_j}\to\infty$.  Normalize their differences from a fixed
$q_0\in Q$ and take a subsequence converging to a vertical unit
vector $(0,v)$.  Closedness and convexity imply
$q_0+t(0,v)\in Q$ for every $t\geq0$, contradicting the bounded
fiber through $q_0$.  This argument also proves that $C$ is closed,
by lifting a convergent sequence in $C$ and extracting a subsequence.
The standard relative-interior identity for linear images of convex
sets gives $\pi_E(\intr Q)=\intr C$.

For each interior $x$, the fiber is bounded and its boundary lies
on $\partial Q$.  Barrier divergence and strict convexity therefore
give a unique interior minimizer $y(x)$.  The equation $F_y=0$
and the implicit-function theorem show that $y(x)$ is $C^2$.
For a direction $h$, set $\widehat h=(h,Dy(x)h)$.  Differentiating
stationarity gives Hessian orthogonality to every vertical direction,
and hence
\begin{align*}
 D\phi[h]&=DF[\widehat h],\\
 D^2\phi[h,h]&=D^2F[\widehat h,\widehat h]
                 =\min_kD^2F[(h,k),(h,k)],\\
 D^3\phi[h,h,h]&=D^3F[\widehat h,\widehat h,\widehat h].
\end{align*}
In the third identity the acceleration term is vertical and its
Hessian pairing with $\widehat h$ vanishes.  The positive definite
Schur complement proves nondegeneracy; the two defining derivative
inequalities for $F$ prove the same inequalities for $\phi$.
Finally, as $x_j$ tends to a finite boundary point of $C$, the
minimizers remain bounded by the first paragraph.  Any cluster point
lies on $\partial Q$, since an interior cluster point would project
to an interior point of $C$.  Barrier divergence therefore implies
$\phi(x_j)\to+\infty$.  This completes the proof.
\end{proof}

Let $\mathbb F\in\{\R,\C,\HH\}$, let $p,b\geq1$,
and choose nonzero real linear subspaces $E_a\subset\mathbb F^p$.
With $W=[w_1\ \cdots\ w_b]$, $w_a\in E_a$, define
\begin{equation}\label{eq:packing-domain-concrete}
 \Omega_{\mathrm{pack}}
 =\left\{(S,W):S\in H^b(\mathbb F),\ S_{aa}=1,\
       \begin{pmatrix}S&W^*\\W&I_p\end{pmatrix}\succ0\right\}.
\end{equation}
Determinants for quaternionic Hermitian matrices mean Jordan
determinants, and traces in the calculations below mean real traces.

\begin{theorem}[Intrinsic column-packing parameter]\label{thm:packing-intrinsic}
For the single-factor domain in \eqref{eq:packing-domain-concrete},
\begin{equation}\label{eq:packing-intrinsic}
 \vartheta_{\mathrm{opt}}(\Omega_{\mathrm{pack}})=b.
\end{equation}
The restricted standard barrier
$F_{\mathrm{pack}}=-\log\det(S-W^*W)$ attains this value.
This holds without any relation between the number of columns $b$
and the row dimension $p$.
\end{theorem}
\begin{proof}
With the Schur complement $D=S-W^*W$, the projection of
$\Omega_{\mathrm{pack}}$ onto $W$ is exactly
$\prod_a\{w_a\in E_a:\norm{w_a}<1\}$.
Indeed its diagonal entries are $d_a=1-\norm{w_a}^2$;
conversely, set $S=W^*W+\operatorname{Diag}(d_a)$.
For fixed $W$, the closed set of feasible $S$ satisfies
$|D_{ac}|^2\leq d_ad_c\leq1$ and is bounded.
The whole closed slice is bounded as well: $\norm{w_a}\leq1$,
and $S\succeq0$ with diagonal one implies $|S_{ac}|\leq1$.
It has relative interior, for example $S=I_b,W=0$.
Lemma~\ref{lem:bounded-fiber-barrier} and the product-of-balls
parameter established above show that every barrier has $\nu\geq b$.

For the upper bound work first on the unrestricted matrix epigraph
$S-W^*W\succ0$.  It is an affine slice of $H^{p+b}_{++}(\mathbb F)$,
so its log-determinant is standard self-concordant.  An affine shear
at any base point sets its $W$ coordinate to zero while preserving
the Schur complement.  At the resulting point $(D,0)$,
\[
 DF[H,K]=-\tr(D^{-1}H),\qquad
 D^2F[(H,K),(H,K)]
 =\tr(D^{-1}HD^{-1}H)+2\tr(D^{-1}K^*K).
\]
The squared dual norm of this gradient is $b$, attained by
$(H,K)=(-D,0)$.  Therefore this is a $b$-self-concordant barrier;
the diagonal and column-subspace restrictions cannot increase its
parameter.  The lower bound proves equality.
\end{proof}

For the standard barrier the partial minimizer is explicit.
Hadamard's determinant inequality, with equality exactly at a diagonal
positive matrix, gives
\[
 \min_{S:\,S_{aa}=1}F_{\mathrm{pack}}(S,W)
 =-\sum_{a=1}^b\log(1-\norm{w_a}^2).
\]
Thus column packing reduces the ambient parameter $p+b$ to $b$ on the
fixed slice, and no barrier there, coupled or not, does better.
In contrast, the optimal parameter of a grouped slice may be much larger
than the number of balls.  The factor count (here one), the ambient
rank, and the optimal slice parameter measure different costs.

\begin{example}[Repeated boundary blocks do not bound arbitrary barriers]
\label{ex:duplicated-barrier}
For $m\geq2$, the diagonal affine slice
\[
 \left\{\bigl((1/m,x/m),\ldots,(1/m,x/m)\bigr):
                          x\in\intr\B_2^s\right\}\subset Q_{s+1}^m
\]
is affinely isomorphic to one ball.  All $m$ blocks reach nonzero
boundary points simultaneously, and polynomial slack factors are
$A_j(x)=(1,x)/m$, $B_j(z)=(1,-z)$.
The restricted standard product barrier is
$-m\log(1-\norm{x}^2)$ plus a constant, with exact parameter $m$;
its squared dual gradient norm approaches $m$ as $\norm{x}\to1$.
Yet the slice admits the one-parameter barrier
$-\log(1-\norm{x}^2)$.  The active conormals coincide after
restriction.  Consequently simultaneous boundary multiplicity
alone does not give a lower bound for arbitrary barriers;
the polyhedral sections or the projection argument above provide the
missing geometric information.
\end{example}

These results concern the parameters of self-concordant barriers on
specified formulations.  They do not assert that every algorithm
takes a number of iterations proportional to the square root of the
parameter, and they do not bound the conditioning or the arithmetic
cost of a Newton solve.
